\documentclass[11pt]{article}
\usepackage[margin=1in]{geometry}
\usepackage{amsmath,amssymb,amsthm,mathtools}
\usepackage[T1]{fontenc}
\usepackage{lmodern,microtype}
\usepackage[colorlinks=true,linkcolor=blue,urlcolor=blue,citecolor=blue]{hyperref}
\newtheorem{theorem}{Theorem}
\newtheorem{lemma}{Lemma}
\newcommand{\F}{\mathcal F}
\newcommand{\N}{\mathcal N}
\newcommand{\Tr}{\operatorname{Tr}}
\newcommand{\supp}{\operatorname{supp}}
\newcommand{\id}{\mathbf 1}
\title{Maximal norm discontinuity for a bounded-region\newline continuum fermion interaction}
\author{Alec Kriebel\\\small ORCID: \href{https://orcid.org/0009-0001-9320-500X}{0009-0001-9320-500X}}
\date{Research note, version 1.0 --- 14 October 2026}
\begin{document}
\maketitle
\begin{abstract}
For continuum fermions in three dimensions, we give one fixed smooth compactly
supported pair potential and one fixed bounded interaction region for which
the Heisenberg orbit of a fixed annihilation operator is discontinuous in
operator norm. Explicit finite-particle Slater states and times tending to zero
witness norm differences tending to the maximal value~$2$. The construction
works both with diagonal density pairs and with normal ordering, without
ultraviolet smearing. A fixed-cutoff commutator estimate, a Slater second-moment
identity, and a two-sector comparison supply a self-contained proof. Continuity
on every fixed vector and in every fixed normal state remains valid.
\end{abstract}

\paragraph{Question and scope.}
Oberwolfach Report 7/2025, printed pp.~365 and~371 \cite{OWR}, distinguishes
pointwise operator-norm continuity from strong-operator continuity and asks
for a potential and states violating the former for an unregularized
interaction restricted to a bounded region. We answer that existential
question in dimension three, with kinetic energy on all of $\mathbb R^3$.
The Hartree discussion preceding the fermionic question concerns a separate
bosonic model. We impose no fixed-density or fixed-energy restriction on the
varying witnesses.

Siebert's earlier 2026 preprint \cite{Discontinuity} proves full-space norm
discontinuity by a high-density cluster and localized phase argument. That
proof adapts directly to the bounded symmetric cutoff kernel, yielding the
qualitative conclusion here. Its Proposition~2.6 also gives the maximal limit~$2$
for the full-space constant-number reference. Our explicit free-Slater
comparison transfers that limit to the fixed spatial cutoff. The audit also
records an additional energy-bootstrap derivation using the earlier localized
phase mechanism; neither cutoff deduction is a theorem literally stated in
the cited preprint. Siebert's 2025 construction \cite{Thermodynamic} uses an
extended algebra, with norm continuity after time averaging. This note claims neither
first priority nor that the 2025 question remained open when this proof was
prepared. The accompanying verification archive records the bounded
primary-source priority audit and the precise relation to earlier mechanisms.

\paragraph{Model and statement.}
Let $\mathfrak h=L^2(\mathbb R^3)$, $h=-\Delta$, and
$\F=\bigoplus_{n\geq0}\bigwedge^n\mathfrak h$. Inner products are linear in
the second variable; exterior products use normalized Slater determinants.
Write $a(f)$ and $a^*(f)$ for the standard CAR operators, with
$\{a(f),a^*(g)\}=\langle f,g\rangle\id$ and $\|a(f)\|=\|f\|$.
The ordering flag $\kappa\in\{0,1\}$ below is distinct from an ultraviolet
regularization parameter; neither model is smeared.

\begin{theorem}\label{thm}
Set $\Lambda=B(0,4)$ and $L=d\Gamma(\id_\Lambda)$. There exist a fixed
real radial $V\in C_c^\infty(\mathbb R^3)$, $0\leq V\leq1$, equal to one
on $\Lambda-\Lambda$, and a fixed normalized
$f\in C_c^\infty(B(0,1/2))$, with the following property. For
\begin{equation}\label{H}
 H^{(\kappa)}_n=T_n+L_n(L_n-\kappa),\qquad
 T_n=\sum_{j=1}^n(-\Delta_{x_j}),\quad
 L_n=\sum_{j=1}^n\id_\Lambda(x_j),
\end{equation}
and their self-adjoint maximal direct sum $H^{(\kappa)}$, there are explicit
unit Slaters $\Psi_N\in\bigwedge^N\mathfrak h$ and
$\Xi_N=a^*(f)\Psi_N\in\bigwedge^{N+1}\mathfrak h$ such that
\begin{equation}\label{times}
 t_N^{(\kappa)}=\frac{\pi}{2N+1-\kappa},\qquad
 D_N^{(\kappa)}=e^{it_N^{(\kappa)}H^{(\kappa)}}a(f)
               e^{-it_N^{(\kappa)}H^{(\kappa)}}-a(f)
\end{equation}
satisfy
\begin{equation}\label{limits}
 \|D_N^{(\kappa)}\Xi_N+2\Psi_N\|\longrightarrow0,\qquad
 \|D_N^{(\kappa)}\Xi_N\|\longrightarrow2,\qquad
 \|D_N^{(\kappa)}\|\longrightarrow2 .
\end{equation}
The same operator-norm limit holds for the creation difference. The normal
vector states $\omega_N(B)=\langle\Xi_N,B\Xi_N\rangle$ satisfy
$\omega_N((D_N^{(\kappa)})^*D_N^{(\kappa)})\to4$.
\end{theorem}

\paragraph{Fixed potential and sector domains.}
Put $b(s)=e^{-1/s}$ for $s>0$ and $b(s)=0$ for $s\leq0$, and choose
\begin{equation}\label{V}
 V(x)=\frac{b(81-|x|^2)}{b(81-|x|^2)+b(|x|^2-64)} .
\end{equation}
The denominator is positive everywhere. This smooth function equals one on
$\overline{B(0,8)}$ and vanishes outside $B(0,9)$. Thus the rigorous
sector realization of the density-density interaction in $\Lambda\times\Lambda$
is
\[
 \sum_{i,j=1}^n\id_\Lambda(x_i)\id_\Lambda(x_j)V(x_i-x_j)=L_n^2.
\]
Omitting the diagonal terms gives $L_n(L_n-1)$, since $V(0)=1$.
We use the coefficient $1$ in this ordered-pair sum, without a factor $1/2$.
These multiplication operators define the field-expression conventions
without multiplying distributions. Each is real, symmetric in the particle
coordinates, and bounded by $n^2$. Hence \eqref{H} is self-adjoint on
$D(T_n)=H^2(\mathbb R^{3n})\cap\bigwedge^n\mathfrak h$. The vacuum sector
is zero. The full domain is
\[
 D(H^{(\kappa)})=\left\{(\eta_n)\in\F:
 \eta_n\in D(T_n),\ \sum_n\|H_n^{(\kappa)}\eta_n\|^2<\infty\right\}.
\]
No uniform Fock-space bound on the interaction is presumed. Nor is invariance
of the CAR algebra under the interacting evolution needed for this norm
comparison in $B(\F)$.

\paragraph{Packed Slaters.}
Choose $f(x)=c_f b(1-9|x|^2)$, with $c_f>0$ normalizing its norm.
Its support is $\overline{B(0,1/3)}$. Let
$Q=2e_1+(-1/4,1/4)^3\subset B(0,3)$, disjoint from $\supp f$, and
normalize
$\varphi(x)=c_\varphi\prod_{k=1}^3 b(1-16(x_k-1/2)^2)$.
This belongs to $C_c^\infty((0,1)^3)$. For $m_N=\lceil N^{1/3}\rceil$,
subdivide $Q$ into $m_N^3$ cubes of side $\ell_N=(2m_N)^{-1}$.
Choose the first $N$ lower corners $c_{j,N}$ in lexicographic order, and set
\[
 \psi_{j,N}(x)=\ell_N^{-3/2}\varphi((x-c_{j,N})/\ell_N),\quad
 \Psi_N=\psi_{1,N}\wedge\cdots\wedge\psi_{N,N},\quad
 \Xi_N=a^*(f)\Psi_N .
\]
Disjoint supports give orthonormal orbitals and $a(f)\Psi_N=0$, so both
Slaters have norm one. Since $m_N\leq2N^{1/3}$,
\[
 \sum_{j=1}^N\|\nabla\psi_{j,N}\|^2
 =N\ell_N^{-2}\|\nabla\varphi\|^2
 \leq16\|\nabla\varphi\|^2N^{5/3}.
\]
For either Slater with its actual particle number $n=N$ or $N+1$, and its
initial orbitals $g_1,\ldots,g_n$, therefore
\begin{equation}\label{energy}
 \sum_{j=1}^n\|g_j\|_{H^1}^2\leq C_0 n^{5/3},\qquad
 C_0=1+16\|\nabla\varphi\|^2+\|\nabla f\|^2 .
\end{equation}
All initial supports lie in $B(0,3)$.

\paragraph{Uniform free leakage.}
Use the single fixed cutoff
\[
 \chi(x)=\frac{b(49/4-|x|^2)}{b(49/4-|x|^2)+b(|x|^2-9)} .
\]
It lies in $C_c^\infty(\Lambda)$ and equals one on $B(0,3)$.
With $q=\id_{\Lambda^c}$,
\begin{align}
 [h,\chi]u&=-(\Delta\chi)u-2\nabla\chi\cdot\nabla u,\nonumber\\
 \|[h,\chi]u\|&\leq C_\chi\|u\|_{H^1},\qquad
 C_\chi=(\|\Delta\chi\|_\infty^2+4\|\nabla\chi\|_\infty^2)^{1/2}.
 \label{comm}
\end{align}
For any initial orbital $g$, differentiate
$e^{-i(s-r)h}\chi e^{-irh}g$ on its $H^2$ domain to obtain
\begin{equation}\label{duhcomm}
 (e^{-ish}\chi-\chi e^{-ish})g
 =-i\int_0^s e^{-i(s-r)h}[h,\chi]e^{-irh}g\,dr .
\end{equation}
The free group preserves $H^2$ and is isometric on $H^1$; multiplication by
$\chi$ preserves $H^2$. Thus the identity is justified for every packed
smooth orbital, while its bound uses only the $H^1$ norm. From $\chi g=g$
and $q\chi=0$, for every real $s$,
\begin{equation}\label{leak}
 \|qe^{-ish}g\|\leq C_\chi|s|\|g\|_{H^1} .
\end{equation}
The indicator $q$ is never differentiated.

\paragraph{Second moment of the leaked number.}
Set $M_n=n-L_n=d\Gamma(q)$ on the $n$-particle sector.
For either initial Slater $\Phi_n$, free evolution
$\Phi_n(s)=e^{-isT_n}\Phi_n$ remains a normalized Slater.
Let $P$ be its rank-$n$ occupied projection and $\lambda=\Tr(Pq)$.
Expanding the determinant's two-particle marginal gives
\begin{equation}\label{moments}
 \langle M_n\rangle=\lambda,\qquad
 \langle M_n^2\rangle=\lambda^2+\lambda-\Tr(PqPq).
\end{equation}
Indeed the diagonal terms are $\sum_j\langle g_j,qg_j\rangle$,
as $q^2=q$; the terms with $j\ne k$ are
$\langle g_j,qg_j\rangle\langle g_k,qg_k\rangle-
|\langle g_j,qg_k\rangle|^2$. Their sum proves \eqref{moments}.
All traces are finite, and $\Tr(PqPq)\geq0$. With
$\beta=C_\chi\sqrt{C_0}$, \eqref{energy} and \eqref{leak} yield
\begin{align}
 \lambda_n(s)&\leq\beta^2s^2n^{5/3},\nonumber\\
 \|M_n\Phi_n(s)\|&\leq\sqrt{\lambda_n(s)+\lambda_n(s)^2}
 \leq\beta|s|n^{5/6}+\beta^2s^2n^{5/3}.\label{Mbound}
\end{align}
This variance estimate concerns Slaters; it is not asserted for general
correlated or interacting states.

\paragraph{Two-sector comparison.}
Introduce only as a reference
$K_n^{(\kappa)}=T_n+n(n-\kappa)$.
Since $M_n$ has spectrum $\{0,\ldots,n\}$,
\begin{equation}\label{difference}
 H_n^{(\kappa)}-K_n^{(\kappa)}
 =-M_n(2n-\kappa-M_n),\qquad
 \|(H_n^{(\kappa)}-K_n^{(\kappa)})\eta\|\leq2n\|M_n\eta\|.
\end{equation}
The difference is bounded on each sector, so strong Duhamel holds there.
Reference evolution differs from free evolution by a scalar phase. Applying
\eqref{Mbound}, for either $\Phi_n$ and every real $t$, gives
\begin{align}
 E_n(t)&:=\|(e^{-itH_n^{(\kappa)}}-e^{-itK_n^{(\kappa)}})\Phi_n\|
 \leq2n\int_0^{|t|}\|M_n\Phi_n(\operatorname{sgn}(t)s)\|\,ds\nonumber\\
 &\leq\beta n^{11/6}|t|^2+\tfrac23\beta^2n^{8/3}|t|^3.
 \label{E}
\end{align}
Thus for \emph{both} $n=N$ and $n=N+1$, since
$n\leq2N$ and $t_N^{(\kappa)}\leq\pi/(2N)$,
\begin{equation}\label{decay}
 E_n(t_N^{(\kappa)})=O(N^{-1/6}+N^{-1/3})\longrightarrow0 .
\end{equation}
The constants are independent of $N$ and $\kappa$. We have used the Slater
estimate only on free reference states, with no assumption that interacting
evolution preserves Slaters or that a changing many-particle free state stays
close to its initial vector.

\paragraph{Phase and norm witness.}
Let $A=a^*(f)$ and, suppressing the fixed ordering flag, write
$F_N^H(t)=\langle\Xi_N,e^{itH}Ae^{-itH}\Psi_N\rangle$.
Replacing the separately evolved bra and ket by their reference evolutions,
unitarity and $\|A\|=1$ imply
\begin{equation}\label{Fbound}
 |F_N^H(t)-F_N^K(t)|\leq E_{N+1}(t)+E_N(t).
\end{equation}
For $K=d\Gamma(h)+\N(\N-\kappa)$, free covariance and the CAR give exactly
\begin{align}
 F_N^K(t)&=e^{it(2N+1-\kappa)}
 \langle a^*(f)\Psi_N,a^*(e^{ith}f)\Psi_N\rangle\nonumber\\
 &=e^{it(2N+1-\kappa)}\langle f,e^{ith}f\rangle .\label{phase}
\end{align}
The last equality uses $a(f)\Psi_N=0$; it does not require $e^{ith}f$ to
remain orthogonal to the occupied orbitals. At $t=t_N$ the scalar phase is
$-1$, and $\langle f,e^{it_Nh}f\rangle\to1$. Consequently
$F_N^H(t_N)\to-1$, whereas $F_N^H(0)=1$. Taking adjoints,
\[
 \langle\Psi_N,D_N\Xi_N\rangle
 =\overline{F_N^H(t_N)-1}\longrightarrow-2 .
\]
The universal upper bound $\|D_N\|\leq2$ now proves both norm limits.
More quantitatively, since $f\in D(h)$,
\[
 2\geq\|D_N\Xi_N\|\geq
 2-E_N(t_N)-E_{N+1}(t_N)-t_N\|hf\|.
\]
Finally
$\|D_N\Xi_N+2\Psi_N\|^2
 \leq8+4\Re\langle\Psi_N,D_N\Xi_N\rangle\to0$.
This proves the first limit in \eqref{limits} and the vector-state conclusion,
and completes the proof of Theorem~\ref{thm}.

\paragraph{Topology and boundaries.}
For every fixed vector $\eta$, any self-adjoint $H$ and bounded $B$ obey
\[
 \|(e^{itH}Be^{-itH}-B)\eta\|
 \leq\|B\|\|(e^{-itH}-1)\eta\|+\|(e^{itH}-1)B\eta\|\to0.
\]
For $B=a(f)$, this also implies
$\Tr(\rho D_t^*D_t)\to0$ for every fixed positive trace-class state
$\rho$ of trace one, by its spectral decomposition and the uniform bound~$4$.
The witnesses in the theorem vary with $N$. The construction asserts an
existential counterexample for the chosen ball and potential. It asserts
neither all potentials or regions, fixed-density thermodynamic behavior,
nor CAR non-invariance. Packing in dimension $d$ gives decay exponents
$-1/2+1/d$ and $-1+2/d$; this estimate does not settle dimensions one or two.

\paragraph{Verification and disclosure.}
The accompanying source archive contains the unchanged original mathematical
packet, independent analytic and source audits, finite algebra/numerical
controls, and publication verification safeguards. Finite controls supplement
the written universal proof; numerical diagnostics are not interval-certified
proofs. Source hashes are consistency fingerprints, not signatures. AI tools
were used extensively in developing, checking, and preparing this work,
including adversarial reviews by separate AI agents. This preprint is
unrefereed, and no independent human peer review is claimed.

\begin{thebibliography}{9}
\bibitem{OWR} Oberwolfach Report 7/2025,
\emph{Mini-Workshop: New Directions in Correlated Quantum Systems},
\href{https://doi.org/10.4171/owr/2025/7}{doi:10.4171/owr/2025/7},
printed pp.~364--365 and~371.
\bibitem{Thermodynamic} O.~Siebert,
\emph{On the thermodynamic limit of interacting fermions in the continuum},
Letters in Mathematical Physics \textbf{115}, 33 (2025),
\href{https://doi.org/10.1007/s11005-025-01919-y}{doi:10.1007/s11005-025-01919-y};
\href{https://arxiv.org/abs/2409.10495v2}{arXiv:2409.10495v2}, 24 February 2025.
\bibitem{Discontinuity} O.~Siebert,
\emph{Discontinuity of Continuum Fermion Dynamics},
\href{https://arxiv.org/abs/2609.34402v1}{arXiv:2609.34402v1},
submitted 28 September 2026.
\end{thebibliography}
\end{document}
